\chapter{Universal Computation, Reductions, and Fixed Points}
\label{ch:universal-reductions}

\section{An effective numbering of programs}
\label{sec:universal-reductions-01}

Normalize every machine description as follows.  A description is
\emph{normalized} iff its tape symbols are the numbers
\(0,\ldots,\abs{\mathcal A_{\mathrm{tape}}}-1\) with the designated
assignments \(\sqcup\mapsto0\), digit \(0\mapsto1\), digit \(1\mapsto2\),
and its states are the numbers \(0,\ldots,\abs{Q}-1\) with the designated
assignments \(q_0\mapsto0\), \(q_{\mathrm{acc}}\mapsto1\),
\(q_{\mathrm{rej}}\mapsto2\); the remaining numbers are assigned
arbitrarily, and the row list is the lexicographic listing, by numbered key
\((q,a)\), of the full transition table over its non-halting domain.  (A
clause fixing an ``order of first appearance'' instead would be circular:
the lexicographic row order presupposes the numbering of states and
symbols.  Without the designated assignments, ``valid code'' below would
not be well defined.)  The following serialization fixes the entire
description, including its alphabets, before assigning it an index.

\begin{definition}[Serialization of a normalized numerical machine]
\label{def:p04-program-serialization}
Let \(n\geq3\) be the number of states and \(b\geq3\) the number of
tape symbols.  Encode the input alphabet by the natural number
\[
 A=\sum_{a=0}^{b-1}\eta_a2^a,\qquad
 \eta_a\in\{0,1\},\quad \eta_0=0,\quad\eta_1=\eta_2=1,
\]
where \(\eta_a=1\) iff symbol number \(a\) belongs to the input
alphabet.  This is a fixed-width bitmap: digits above position \(b-1\)
must be zero.  The powers in this numerical coding definition are
metatheoretic operations on natural numbers; no arithmetic-language
exponentiation symbol is introduced.

There are \(m=(n-2)b\) nonhalting state/symbol keys.  Order them
lexicographically, with states \(0,3,4,\ldots,n-1\) and symbols
\(0,\ldots,b-1\).  For the row
\(\delta(q,a)=(q',a',d)\), use the code
\[
 \rho=\operatorname{code}_5(q,a,q',a',\varepsilon),\qquad
 \varepsilon=0\text{ for }d=L,\quad
 \varepsilon=1\text{ for }d=R.
\]
Thus each row has exactly five fields, with
\(q,q'<n\), \(a,a'<b\), \(q\notin\{1,2\}\), and
\(\varepsilon\in\{0,1\}\).  Write the ordered row codes as
\(\rho_0,\ldots,\rho_{m-1}\).  Each allowed key occurs exactly
once, and no row has a halting source state.

The description tuple and its index are exactly
\[
 \bar u=(n,b,A,\rho_0,\ldots,\rho_{m-1}),\qquad
 k=3+m,\qquad
 e=\operatorname{pair}_{\N}
       \bigl(k,\operatorname{code}_k(\bar u)\bigr).
\]
The tuple code is the previously defined right-nested Cantor code, without
an additional terminal marker.  A number is a valid program code iff
decoding gives precisely this length, bitmap and full ordered row list,
with all the displayed ranges and designated-slot conditions satisfied.
\index{program!serialization}\index{machine code!validity}
\end{definition}

Decoding is unique and effective.  Invert the outer pairing to obtain
\((k,c_0)\), then decode the right-nested tuple of the specified length:
for \(k>1\), successively invert
\(c_i=\operatorname{pair}_{\N}(u_i,c_{i+1})\) for \(i<k-1\),
and put \(u_{k-1}=c_{k-1}\); the one-entry code is the entry itself,
and a zero-entry tuple has code \(0\).  Tuples of length less than
\(6\) are invalid machine descriptions because \(n,b\geq3\)
forces \(k=3+(n-2)b\geq6\).  For the remaining tuples, read the
first three fields as \(n,b,A\), check the required length and bitmap
by finite binary division, decode each row by four pairing inversions,
and check its range and its place in the full lexicographic key list.
Every loop is finite; the effective inverse of the pairing map and
finite integer arithmetic therefore decide validity.  Conversely, every
normalized numerical machine supplies exactly such metadata and rows,
and decoding its code recovers its full transition function and input
alphabet.  This verifies the encoding/decoding interface used below.

No canonical machine-to-code map is fixed: the permitted renamings of
nondesignated states and symbols can give different indices, and indices
need not be unique.  A natural number that is not the code of a
valid normalized numerical-machine description denotes a machine that
diverges on every input.  Such a description must in particular satisfy
\(\{0,1\}\subseteq\mathcal A_{\mathrm{in}}\), read at the designated digit
numbers \(1,2\) together with the input and output conventions fixed
above; a normalized description whose input alphabet omits either binary
symbol is treated as invalid by this numbering.  In a valid code the three
designated states are pairwise distinct; a machine whose start state is
halting, which the machine definition permits by requiring only
\(q_{\mathrm{acc}}\neq q_{\mathrm{rej}}\), is captured up to behavior by
the surrogate-start step in the proof of \cref{thm:enumeration} below.
For a valid code \(e\), define
\(\varphi_e(x)\) to be the numerical output when that machine accepts with a
valid numerical output, and leave \(\varphi_e(x)\) undefined when it rejects,
diverges, or halts without a valid numerical output.  A fixed wrapper simulates the decoded
machine, copies a canonical numerical output after a valid accepting halt, and
enters an infinite loop after either rejection or an accepting halt without a
valid numerical output.

The number \(e\) is an \emph{index} for \(\varphi_e\).  Different indices may
compute the same partial function.  For a finite tuple \(\bar x\), the notation
\(\varphi_e(\bar x)\), with \(\bar x\in\N^k\), denotes the behavior of indexed
program \(e\) on the input word
\(\operatorname{bin}(\operatorname{code}_k(\bar x))\).  This is
arity-dependent application notation for the one fixed program numbering, not
a second numbering and not a second application of the unary tuple code.

For expressions denoting possibly undefined values, \(a\simeq b\) means that
both sides are undefined or both are defined and equal.  For partial
functions \(f,g:\N\rightharpoonup\N\), the notation \(f\simeq g\)
abbreviates: \(f(x)\simeq g(x)\) for every \(x\in\N\).

\begin{lemma}[Renaming invariance]
\label{lem:renaming-invariance}
Let \(\mathcal T\) be a Turing machine used to compute a numerical function
and let \(\iota=(\iota_Q,\iota_A)\) bijectively rename its states and tape
symbols, preserving the roles of
\(q_0,q_{\mathrm{acc}},q_{\mathrm{rej}}\) and of \(\sqcup\) and the digits
\(0,1\).  Write \(\mathcal T^{\iota}\) for the machine obtained by applying
\(\iota\) to every component of the machine tuple, and \(C^{\iota}\) for the
corresponding renaming of a configuration \(C\).  Then
\(C\longrightarrow_{\mathcal T}C'\) iff
\(C^{\iota}\longrightarrow_{\mathcal T^{\iota}}(C')^{\iota}\).  Hence, on
every pair of input words corresponding letter by letter under \(\iota_A\),
the machines \(\mathcal T\) and \(\mathcal T^{\iota}\) accept, reject, or
diverge together, with corresponding tapes at every step.
\textit{Status: complete proof.}
\end{lemma}

\begin{proof}
For the one-step claim, read the configuration definition on both sides:
\(\delta(q,T(h))=(q',a,d)\) holds in \(\mathcal T\) iff
\(\delta^{\iota}(\iota_Q(q),\iota_A(T(h)))
=(\iota_Q(q'),\iota_A(a),d)\) holds in \(\mathcal T^{\iota}\), since
\(\delta^{\iota}\) is by construction the \(\iota\)-conjugate of \(\delta\);
the rewritten cell and the moved head then correspond position by position.
The claim for whole computations follows by induction on the number of
steps.  Initial configurations correspond because \(\iota_A\) preserves
\(\sqcup\), and halting, acceptance, and rejection correspond because
\(\iota_Q\) preserves the designated states.  Valid numerical outputs
correspond because \(\iota_A\) preserves the digits.
\end{proof}

\begin{theorem}[Adequacy of the numbering: enumeration theorem]
\label{thm:enumeration}
\begin{enumerate}[label=(\alph*)]
\item Every \(\varphi_e\) is partial computable.
\item For every \(k\geq1\) and every partial computable
      \(f:\N^k\rightharpoonup\N\) there is an index \(e\) with
      \(\varphi_e(\bar x)\simeq f(\bar x)\) for all \(\bar x\in\N^k\),
      under the arity-\(k\) application convention fixed above.
\end{enumerate}
\textit{Status: complete proof.}
\end{theorem}

\begin{proof}
(a) For a valid code \(e\), the fixed wrapper described in the numbering
paragraph simulates the decoded machine, accepts with the canonical output
word exactly when that machine accepts with a valid numerical output, and
otherwise runs forever; for an invalid code, \(\varphi_e\) is the
everywhere-undefined function, which a machine with an immediate infinite
loop computes.  In both cases \(\varphi_e\) satisfies
\cref{def:partial-computable}.

(b) Let \(\mathcal T\) witness that \(f\) is partial computable in the sense
of \cref{def:partial-computable}.  First,
\(\{0,1\}\subseteq\mathcal A_{\mathrm{in}}\) holds automatically:
\(\operatorname{code}_1\) is the identity and, for \(k\geq2\),
\(\operatorname{code}_k\) is built from the pairing bijection, so
\(\operatorname{code}_k\) is surjective for every \(k\geq1\) and the input
words \(\operatorname{bin}(\operatorname{code}_k(\bar x))\) exhaust all
binary numerals.  An initial configuration exists only for a word over
\(\mathcal A_{\mathrm{in}}\), and \cref{def:partial-computable} demands an
accepting or infinite computation of \(\mathcal T\) on each of these words;
since some required numeral contains both digits, both digits lie in
\(\mathcal A_{\mathrm{in}}\), and no alphabet extension is needed.

Second, since the machine definition requires only
\(q_{\mathrm{acc}}\neq q_{\mathrm{rej}}\), the start state of
\(\mathcal T\) may coincide with a halting state.  In that case replace it
by a fresh start state \(q_0'\) with \(\delta(q_0',a)=(q_0,a,R)\) for every
tape symbol \(a\): the one extra step leaves the tape unchanged, and the
output convention ignores the final head position and the cells outside the
output prefix, so the accepted outputs and the divergence behavior are
unchanged.  After this step the three designated states are pairwise
distinct, as are the three designated symbols \(\sqcup,0,1\).

Now choose any bijective renaming \(\iota\) of states and tape symbols onto
the designated-slot numbering that preserves the six designated roles, and
let \(\mathcal T^*\) be the resulting normalized machine;
\cref{lem:renaming-invariance} makes \(\mathcal T^*\) behave exactly as
\(\mathcal T\) on corresponding inputs, and the binary numerals correspond
to themselves read at the designated digit numbers \(1,2\).  List the rows
of \(\mathcal T^*\) lexicographically by numbered key \((q,a)\); the row
list is the full graph of the transition function, which is total on
non-halting pairs by the machine definition.  Encode the resulting
description tuple of Definition~\ref{def:p04-program-serialization},
including its state count, tape-alphabet size and input-alphabet bitmap,
with length \(k'\), as
\(\operatorname{pair}_{\N}(k',\operatorname{code}_{k'}(\bar n))\), obtaining
a number \(e\).  This \(e\) is a valid code by construction.  On each input
word \(\operatorname{bin}(\operatorname{code}_k(\bar x))\), the decoded
machine \(\mathcal T^*\) accepts with output word
\(\operatorname{bin}(f(\bar x))\) when \(f(\bar x)\) is defined and runs
forever when \(f(\bar x)\) is undefined, so
\(\varphi_e(\bar x)\simeq f(\bar x)\) for all \(\bar x\in\N^k\).
\end{proof}

Kleene's treatment of enumeration and normal form over standard numberings
\cite{Kleene1952} is parallel reading; the in-document proof is primary
because this numbering is bespoke: its normalization, its validity
constraint \(\{0,1\}\subseteq\mathcal A_{\mathrm{in}}\), and its wrapper are
fixed here rather than taken from the cited source.

\begin{theorem}[Universal computation]
\label{thm:universal-computation}
There is a partial computable function \(U(e,x)\) such that
\[
  U(e,x)\simeq\varphi_e(x)
\]
for all \(e,x\in\N\).
\textit{Status: proof sketch.}
\end{theorem}

\begin{proofplan}
Decode \(e\) as a machine description.  If decoding fails, diverge.  Otherwise
store an encoded configuration and repeatedly apply the machine's finite
 transition table.  Return an output only when the simulated machine reaches
its accepting state with a valid numerical output; copy that output into
canonical form.  If it rejects or accepts without a valid numerical output,
enter an
 infinite loop.  Each simulation step is effective, and every finite computation is
eventually reproduced.  Turing's original universal-machine construction is
the historical source \cite{Turing1937}.
\end{proofplan}

\begin{theorem}[Parameter theorem]
\label{thm:parameter}
For every \(m,n\geq1\), there is a total computable function
\(\mathsf{par}_n^m:\N^{m+1}\to\N\) such that, for every index \(e\), every \(\bar a\in\N^m\), and every
\(\bar x\in\N^n\),
\[
  \varphi_{\mathsf{par}_n^m(e,\bar a)}(\bar x)
  \simeq \varphi_e(\bar a,\bar x).
\]
\textit{Status: cited theorem.}
\end{theorem}

From \(e\) and the constants \(\bar a\), construct a normalized transition
table for a machine that decodes its runtime input as the \(n\)-tuple
\(\bar x\), computes
\(\operatorname{code}_{m+n}(\bar a,\bar x)\), writes the binary word for that code,
and then simulates program \(e\).  The established tuple encoding and decoding
maps, state renumbering, and finite table concatenation are computable, so the
map from \((e,\bar a)\) to the new index is total and computable.  The omitted
verification writes the compiler's transition table and proves its simulation
invariant.  See \cite{Kleene1952} for a complete proof of the parameter
theorem for effective program numberings.
